% Part: sets-functions-relations
% Chapter: relations
% Section: orders

\documentclass[../../../include/open-logic-section]{subfiles}

\begin{document}

\olfileid{sfr}{rel}{ord}
\olsection{Ordeningen}

\begin{explain}
Veel vergelijkingen drukken uit dat objecten in een bepaald opzicht
``kleiner dan'', ``gelijk aan'' of ``groter dan'' andere objecten zijn. Daarbij
gaat het om \emph{ordeningsrelaties}, waarvan verschillende soorten bestaan.
Soms moet elk tweetal objecten vergelijkbaar zijn, soms niet. Sommige relaties
bevatten de identiteit (zoals~$\le$), andere niet (zoals~$<$). Een indeling is
daarom nuttig.
\end{explain}

\begin{defn}[Preorde]
Een relatie die zowel reflexief als transitief is, heet een
\emph{preorde.}  
\end{defn}

\begin{defn}[Partiële ordening]
Een preorde die bovendien antisymmetrisch is, heet een
\emph{partiële ordening}.
\end{defn}

\begin{defn}[Lineaire ordening]\ollabel{def:linearorder}
Een partiële ordening die bovendien samenhangend is, heet een
\emph{totale ordening} of \emph{lineaire ordening}.
\end{defn}

\begin{ex}
Elke lineaire ordening is ook een partiële ordening en elke partiële ordening is
ook een preorde, maar de omkeringen gelden niet. De universele relatie
op~$A$ is een preorde, want zij is reflexief en transitief. Als
$A$ meer dan één !!{element} bevat, is de universele relatie echter niet
antisymmetrisch en dus geen partiële ordening.
\end{ex}

\begin{ex}
Beschouw de relatie \emph{niet langer dan} $\preccurlyeq$ op~$\Bin^*$:
$x \preccurlyeq y$ dan en slechts dan als $\len{x} \le \len{y}$. Dit is een
preorde (reflexief en transitief), en zelfs samenhangend, maar geen partiële
ordening, want de relatie is niet antisymmetrisch. Zo geldt $01
\preccurlyeq 10$ en $10 \preccurlyeq 01$, terwijl $01 \neq 10$.
\end{ex}

\begin{ex}
Een belangrijke partiële ordening is de relatie $\subseteq$ op een verzameling
van verzamelingen. In het algemeen is dit geen lineaire ordening: als $a \neq b$
en we $\Pow{\{a, b\}} = \{\emptyset, \{a\}, \{b\}, \{a,b\}\}$ beschouwen,
zien we dat $\{a\} \nsubseteq \{b\}$, $\{a\} \neq \{b\}$ en $\{b\}
\nsubseteq \{a\}$.
\end{ex}

\begin{ex}
De relatie \emph{deelbaarheid zonder rest} definieert een partiële ordening die
niet lineair is. Voor gehele getallen $n$ en~$m$ betekent $n \mid m$ dat $n$
het getal $m$ zonder rest deelt, wat precies het geval is als er een geheel
getal~$k$ bestaat met $m = kn$. Op~$\Nat$ is dit een partiële
ordening, maar geen lineaire ordening: zo geldt $2 \nmid 3$ en ook $3 \nmid
2$. Als relatie op $\Int$ is deelbaarheid slechts een preorde, want zij is niet
antisymmetrisch: $1 \mid -1$ en $-1 \mid 1$, maar $1 \neq -1$.
\end{ex}

\begin{ex}
De \emph{uitbreidingsrelatie} op de verzameling rijen~$A^*$ wordt als volgt
gedefinieerd: $s \sqsubseteq s'$ dan en slechts dan als $s = \emptyseq$ (de lege
rij), $s = s'$, of $s = \tuple{s_1, \dots, s_n}$ en $s' =
\tuple{s_1, \dots, s_n, s_{n+1}, \dots, s_m}$. Als $s \sqsubseteq s'$,
zeggen we ook dat $s$ een \emph{beginsegment} van~$s'$ is. De
uitbreidingsrelatie op $A^*$ is een partiële maar geen lineaire ordening; als
bijvoorbeeld $a \neq b$, dan geldt $ab \not\sqsubseteq ba$ en $ba \not\sqsubseteq ab$.
\end{ex}

\begin{defn}[Strikte ordening]
Een \emph{strikte ordening} is een relatie die irreflexief, asymmetrisch en
transitief is.
\end{defn}

\begin{defn}[Strikte lineaire ordening]\ollabel{def:strictlinearorder}
Een strikte ordening die bovendien samenhangend is, heet een
\emph{strikte totale ordening} of \emph{strikte lineaire ordening}.
\end{defn}

\begin{ex}
De lineaire ordening~$\le$ hoort bij de strikte lineaire ordening~$<$; de
partiële ordening~$\subseteq$ bij de strikte ordening~$\subsetneq$.
\end{ex}

Door aan een strikte ordening $R$ op~$A$ de diagonaal $\Id{A}$ toe te voegen,
dus alle paren~$\tuple{x, x}$, ontstaat een partiële ordening.  (Dit heet de
\emph{reflexieve afsluiting} van~$R$.) Omgekeerd ontstaat uit een partiële
ordening een strikte ordening door~$\Id{A}$ te verwijderen. De volgende twee
resultaten formuleren dit precies.

\begin{prop}\ollabel{prop:stricttopartial}
Als $R$ een strikte ordening op~$A$ is, dan is $R^+ = R \cup \Id{A}$ een
partiële ordening. Is $R$ bovendien een strikte lineaire ordening, dan is
$R^+$ een lineaire ordening.
\end{prop}

\begin{proof}
Stel dat $R$ een strikte ordening is, dus dat $R \subseteq A^2$ en $R$
irreflexief, asymmetrisch en transitief is. Zij $R^+ = R \cup \Id{A}$. We
moeten aantonen dat $R^+$ reflexief, antisymmetrisch en transitief is.

$R^+$ is duidelijk reflexief, want $\tuple{x, x} \in \Id{A} \subseteq
R^+$ voor alle $x \in A$. 

Voor de antisymmetrie van $R^+$ nemen we ter tegenspraak aan dat
$R^+xy$ en $R^+yx$, terwijl $x \neq y$. Nu is $\tuple{x,y} \in R \cup \Id{A}$,
maar $\tuple{x, y} \notin \Id{A}$, zodat $\tuple{x, y} \in R$, oftewel
$Rxy$. Evenzo geldt~$Ryx$. Dit is echter in tegenspraak met de aanname
dat $R$ asymmetrisch is.

Voor de transitiviteit nemen we aan dat $R^+xy$ en $R^+yz$. Als zowel
$\tuple{x, y} \in R$ als $\tuple{y,z} \in R$, dan is $\tuple{x, z} \in
R$, omdat $R$ transitief is. Anders geldt $\tuple{x, y} \in
\Id{A}$, dus $x = y$, of $\tuple{y, z} \in \Id{A}$, dus $y = z$.
In het eerste geval geldt volgens de aanname $R^+yz$ en volgt uit $x = y$ dat
$R^+xz$. Voor het tweede geval geldt hetzelfde. In beide gevallen is $R^+xz$;
dus is $R^+$ ook transitief.

Neem ten slotte aan dat $R$ ook samenhangend is.
Voor alle $x \neq y$ geldt dan $Rxy$ of~$Ryx$, oftewel
$\tuple{x, y} \in R$ of $\tuple{y, x} \in R$. Omdat $R \subseteq R^+$,
blijft dit voor $R^+$ gelden, zodat ook $R^+$ samenhangend is.
\end{proof}

\begin{prop}\ollabel{prop:partialtostrict}
Als $R$ een partiële ordening op~$A$ is, dan is $R^- = R \setminus \Id{A}$ een
strikte ordening. Is $R$ bovendien een lineaire ordening, dan is $R^-$ een
strikte lineaire ordening.
\end{prop}

\begin{proof}
Het bewijs wordt als opgave overgelaten.
\end{proof}

\begin{prob}
Geef een bewijs van \olref[sfr][rel][ord]{prop:partialtostrict}. 
\end{prob}

Het volgende eenvoudige resultaat laat zien dat strikte lineaire ordeningen
een op extensionaliteit lijkende eigenschap hebben:

\begin{prop}\ollabel{prop:extensionality-strictlinearorders}
Als $<$ een strikte lineaire ordening op $A$ is, dan geldt:
\[
  (\forall a, b \in A)((\forall x \in A)(x < a \liff x < b) \lif a = b).
\]
\end{prop}

\begin{proof}
Stel $(\forall x \in A)(x < a \liff x < b)$. Als $a < b$, dan $a <
a$, in tegenspraak met de irreflexiviteit van $<$; dus $a \nless b$.
Geheel analoog geldt $b \nless a$. Dan volgt $a = b$ uit de samenhangendheid
van~$<$.
\end{proof}
\end{document}
